\documentclass[10pt,aspectratio=169,t]{beamer}

\usepackage[T1]{fontenc}
\usefonttheme{serif}
\usepackage{mathpazo}
\usepackage{booktabs}
\usepackage{array}
\usepackage{listings}
\usepackage{xcolor}

% ---------- theme: clean, compact (same as findings_AB.tex) ----------
\definecolor{accent}{HTML}{1F4E79}
\usetheme{default}
\setbeamertemplate{navigation symbols}{}
\setbeamercolor{structure}{fg=accent}
\setbeamercolor{frametitle}{fg=accent}
\setbeamercolor{title}{fg=accent}
\setbeamercolor{itemize item}{fg=accent}
\setbeamercolor{itemize subitem}{fg=accent}
\setbeamerfont{frametitle}{size=\large,series=\bfseries}
\setbeamerfont{title}{size=\LARGE,series=\bfseries}
\setbeamerfont{subtitle}{size=\normalsize}
\setbeamersize{text margin left=1.4em,text margin right=1.4em}
\setbeamertemplate{frametitle}{%
  \vskip1.2ex
  \usebeamerfont{frametitle}\usebeamercolor[fg]{frametitle}\insertframetitle\par
  \vskip0.4ex
  {\color{accent}\hrule height 0.6pt}
  \vskip0.6ex}
\setbeamertemplate{footline}{%
  \hfill{\scriptsize\color{gray}\insertframenumber\,/\,\inserttotalframenumber}\hspace{1.4em}\vskip1.5ex}
\setbeamertemplate{itemize items}[circle]
\setlength{\leftmargini}{1.2em}

\lstdefinestyle{code}{
  basicstyle=\ttfamily\scriptsize,
  columns=fullflexible,
  keepspaces=true,
  showstringspaces=false,
  frame=single,
  rulecolor=\color{gray!60},
  backgroundcolor=\color{gray!6},
  xleftmargin=2pt,
  xrightmargin=2pt,
  aboveskip=3pt,
  belowskip=3pt,
  breaklines=true,
  breakatwhitespace=true,
}
\lstdefinestyle{cpp}{style=code, language=C++,
  keywordstyle=\color{blue!60!black}, commentstyle=\color{green!40!black}}
\lstdefinestyle{out}{style=code, backgroundcolor=\color{white}}
\lstdefinestyle{cppsmall}{style=cpp, basicstyle=\ttfamily\tiny}

\newcommand{\spec}[1]{\textcolor{blue!50!black}{\small\itshape #1}}

\title{Weekly Update}
\subtitle{Findings C and D}
\author{}
\date{\today}

\begin{document}

\frame{\titlepage}

% ============================================================
\section{Finding C}

\begin{frame}{Finding C: Handshake on an Established Session}
\begin{center}
\Large Finding C\\[4pt]
\normalsize a second handshake frame arrives while the session is already connected\\[4pt]
\small 4.19.3.1, 4.19.4.3, 4.19.4.6, 4.19.4.11 (BTP) \quad and \quad 4.20.3.1, .3, .6, .11 (PAFTP)\\[2pt]
\footnotesize Page numbers in this deck are printed spec pages (PDF page $-$ 4).
\end{center}
\end{frame}

% ------------------------------------------------------------
\begin{frame}{C --- What the Spec Says}
\small
\begin{itemize}
  \item \textbf{4.19.3.1 (pp.~225--226):} the handshake frame has its own
        layout: flags \texttt{0x65}, opcode \texttt{0x6C}, versions, MTU, window
        size. It has \emph{no sequence-number field}.
  \item \textbf{4.19.4.3 (pp.~228--230):} handshake rules only describe
        establishing a session: \spec{``The handshake request packet SHALL
        include, a list of BTP protocol versions supported by the client
        \ldots''}
  \item \textbf{4.19.4.6 (pp.~231--232):} \spec{``Peers SHALL check to ensure
        that all received BTP packets properly increment the sender's previous
        sequence number by 1. If this check fails, the peer SHALL close the BTP
        session and report an error to the application.''}
  \item \textbf{Figures 26/27 (4.19.4.11, pp.~235--238):} no transition for a
        handshake received in the Connected state (checked by hand).
\end{itemize}
\footnotesize PAFTP: same structure and wording in 4.20.3.1, .3, .6 and
Figures 33/34 (4.20.3.11).
\end{frame}

% ------------------------------------------------------------
\begin{frame}{C --- Why It Is a Candidate Underspecification}
\small
Three spec facts, taken together:
\begin{enumerate}
  \item the handshake frame carries no sequence number (4.19.3.1)
  \item every received packet must pass the sequence check, or the session is
        closed (4.19.4.6)
  \item handshake behaviour is only described for session establishment
        (4.19.4.3)
\end{enumerate}
\vspace{2pt}
Is a second handshake frame on an established session a ``packet'' subject to
the check? The text does not say:
\begin{center}
\small
\begin{tabular}{p{0.38\textwidth}p{0.45\textwidth}}
\toprule
\textbf{Reading 1} & \textbf{Reading 2} \\
\midrule
it is a packet; the check fails; \textbf{close the session}
 & it is a handshake, not data; \textbf{ignore it} \\
\bottomrule
\end{tabular}
\end{center}
\vspace{2pt}
\textbf{Origin:} the Week~12 TLC model flagged a handshake response arriving in
\texttt{Established}. The manual review kept this half and dropped the ``stale
timeout'' half, because the figures show no such transition.
\end{frame}

% ------------------------------------------------------------
\begin{frame}[fragile]{C --- SDK, Bluetooth: No Handshake Flag}
\begin{lstlisting}[style=cppsmall]
// src/ble/BtpEngine.h:81-87  -- the only header flags BTP knows
enum class HeaderFlags : uint8_t {
    kStartMessage = 0x01, kContinueMessage = 0x02,
    kEndMessage   = 0x04, kFragmentAck     = 0x08 };       // no handshake (H) bit

// src/ble/BLEEndPoint.cpp:1143-1174  -- first packet only: a one-shot latch, not a state check
if (!mConnStateFlags.Has(ConnectionStateFlag::kCapabilitiesMsgReceived))
{   ...  ExitNow();  }                                      // handshake handled, not fed to reassembler

// src/ble/BLEEndPoint.cpp:1181 and :1193  -- every later frame is treated as data
if (!IsConnected(mState)) { ... }
err = mBtpEngine.HandleCharacteristicReceived(std::move(data), receivedAck, didReceiveAck);
\end{lstlisting}
\begin{itemize}\small
  \item The latch is set once and never cleared, so a second handshake cannot
        reach the negotiation code
  \item The frame is not recognised as a handshake: it goes down the data path
\end{itemize}
\end{frame}

% ------------------------------------------------------------
\begin{frame}[fragile]{C --- SDK, Bluetooth: the Frame Is Read as Data $\Rightarrow$ Close}
Stray handshake bytes: \texttt{0x65} (flags), \texttt{0x6C} (opcode), versions \ldots
\begin{lstlisting}[style=cppsmall]
// src/ble/BtpEngine.cpp:243-266  (abridged)
err = reader.Read8(rx_flags.RawStorage()).StatusCode();   // byte 0 = 0x65, ack bit (0x08) clear
...
err = reader.Read8(&mRxNewestUnackedSeqNum).StatusCode(); // byte 1 = 0x6C = 108, read as sequence number
VerifyOrExit(mRxNewestUnackedSeqNum == mRxNextSeqNum,     // :263
             err = BLE_ERROR_INVALID_BTP_SEQUENCE_NUMBER);
\end{lstlisting}
\begin{lstlisting}[style=cppsmall]
// src/ble/BLEEndPoint.cpp:1296  -- exit path of Receive()
if (err != CHIP_NO_ERROR) { DoClose(closeFlags, err); }
\end{lstlisting}
\begin{itemize}\small
  \item Unless the expected sequence number happens to be 108, the check fails
        and the session is closed, as 4.19.4.6 requires
  \item If it is 108, the frame is accepted as a start+end data fragment and a
        short message is handed up; lengths are bounds-checked
        (\texttt{BtpEngine.cpp}) and the connected peer can send data anyway
\end{itemize}
\end{frame}

% ------------------------------------------------------------
\begin{frame}[fragile]{C --- SDK, Wi-Fi PAF: the Frame Is Recognised and Dropped}
\begin{lstlisting}[style=cppsmall]
// src/wifipaf/WiFiPAFTP.h:90-91
kManagementOpcode = 0x20,
kHankshake        = 0x40,                  // (sic) handshake bit

// src/wifipaf/WiFiPAFEndPoint.cpp:1046-1050  -- after the first-packet gate
if (rx_flags.Has(WiFiPAFTP::HeaderFlags::kHankshake))
{
    ChipLogDebugWiFiPAFEndPoint(WiFiPAF, "Unexpected handshake packet => drop");
    ExitNow();                             // err == CHIP_NO_ERROR: no close
}
\end{lstlisting}
\begin{itemize}\small
  \item The session survives
  \item The log string states a behaviour; it cites no spec section and does
        not mention the spec being silent
\end{itemize}
\end{frame}

% ------------------------------------------------------------
\begin{frame}{C --- The Two Transports Resolve It Differently}
\begin{center}
\small
\begin{tabular}{p{0.30\textwidth}p{0.28\textwidth}p{0.30\textwidth}}
\toprule
 & \textbf{BTP} (\texttt{src/ble}) & \textbf{PAFTP} (\texttt{src/wifipaf}) \\
\midrule
Handshake bit in header flags & no & yes (\texttt{0x40}) \\
Stray handshake, session established & session closed & frame dropped, session kept \\
Code origin & imported 2020, \texttt{362e2a3fae} & \texttt{f1de3aaaca}, 2025-04-01 (part of the original file) \\
Comment or commit message giving a reason & none & none \\
\bottomrule
\end{tabular}
\end{center}
\begin{itemize}\small
  \item Same spec sentence, same bytes on the wire, opposite behaviour
  \item \texttt{git log -S} finds the PAFTP check only in \texttt{f1de3aaaca};
        a handshake-flag test has never existed in \texttt{BLEEndPoint.cpp}
\end{itemize}
\end{frame}

% ------------------------------------------------------------
\begin{frame}{C --- Ideas We Tested and Dropped}
\small
\begin{itemize}
  \item \textbf{``A second handshake renegotiates the window''} (parameter
        downgrade): \texttt{kCapabilitiesMsgReceived} is set only in the
        first-packet branch (\texttt{BLEEndPoint.cpp:1143--1158}) and never
        cleared. A second handshake cannot reach the negotiation code in either
        transport.
  \item \textbf{``Re-handshake resets sequence numbers, so the honest peer's
        next ack becomes invalid''}: depended on the previous idea.
  \item \textbf{``Stale handshake timeout after connection''}: the figures have
        no such transition (Week~13).
\end{itemize}
\end{frame}

% ------------------------------------------------------------
\begin{frame}{C --- Security Severity}
\textbf{What an attacker needs}
\begin{itemize}\small
  \item A write on characteristic C1 of an \emph{already established}
        connection: the connected peer itself, or injection into that link
  \item Not reachable before a connection exists
\end{itemize}
\textbf{What it gives}
\begin{itemize}\small
  \item BTP: that one session is closed. PAFTP: nothing.
  \item No new capability: the connected peer can already close the session
        (4.19.4.10), and a wrong sequence number or an invalid ack already
        forces a close by spec (4.19.4.6, 4.19.4.8)
  \item No data exposure, no PASE bypass seen
\end{itemize}
\textbf{Assessment:} low severity.\\[2pt]
\small Not verified: nothing was executed; over-the-air injection into an
established connection was not assessed.
\end{frame}

% ============================================================
\section{Related: Window Size}

\begin{frame}{Related to C: the Window Size in the Handshake}
\begin{center}
\Large Window size\\[4pt]
\normalsize how small may the window in the handshake be?\\[4pt]
\small 4.19.4.3, 4.19.4.7 (BTP) \quad and \quad 4.20.3.3, .7 (PAFTP)
\end{center}
\end{frame}

% ------------------------------------------------------------
\begin{frame}{Window --- What the Spec Says, and What Developers Did}
\small
\begin{itemize}
  \item \textbf{4.19.4.7 (pp.~232--233):} \spec{``the largest maximum window size
        any peer may support is 255''}. \textbf{No minimum} is stated; the
        constants in 4.19.5 and 4.20.4 define none.
  \item \textbf{4.19.4.3:} \spec{``The server SHALL select a window size equal
        to the minimum of its and the client's maximum window sizes.''} No rule
        for a window of 0 or 1.
  \item \textbf{4.19.4.7:} a server initialises its counter at
        \spec{``(negotiated maximum window size - 1)''}, which is $-1$ for a
        window of 0.
\end{itemize}
\vspace{2pt}
\begin{center}
\footnotesize
\begin{tabular}{lll}
\toprule
\textbf{Date} & \textbf{Commit} & \textbf{Change} \\
\midrule
2026-02-16 & \texttt{c3480ff4d4} & PAFTP, both sides: reject window $<3$ (``underflow'') \\
2026-06-10 & \texttt{28e77bdbc9} & BTP central side: reject 0 (``DoS'', ``underflow'') \\
2026-07-17 & \texttt{5b4abe1fe1} & BTP peripheral side: reject 0 \\
\bottomrule
\end{tabular}
\end{center}
\small Residual difference: PAFTP minimum is 3, BTP only requires $>0$. The spec
text is unchanged.
\end{frame}

% ------------------------------------------------------------
\begin{frame}[fragile]{Window --- SDK and Consequence (by Reading)}
\begin{lstlisting}[style=cppsmall]
// src/ble/BLEEndPoint.cpp:1014-1016  (peripheral side; first line added by 5b4abe1fe1)
VerifyOrReturnError(req.mWindowSize > 0, CHIP_ERROR_INVALID_ARGUMENT);
mRemoteReceiveWindowSize = mLocalReceiveWindowSize = mReceiveWindowMaxSize =
    std::min(req.mWindowSize, static_cast<uint8_t>(BLE_MAX_RECEIVE_WINDOW_SIZE));

// src/ble/BLEEndPoint.cpp:224  (HandleSubscribeReceived): one-byte counter, window 0 -> 255
mRemoteReceiveWindowSize = static_cast<SequenceNumber_t>(mRemoteReceiveWindowSize - 1);

// src/wifipaf/WiFiPAFConfig.h:107
#define PAF_MIN_RECEIVE_WINDOW_SIZE 3
\end{lstlisting}
\begin{itemize}\small
  \item Before the fix, window 0: the counter wraps to 255; after the first ack
        the sender believes the receiver's window is closed and stops sending
  \item Window 1: data cannot be sent, only acks (\texttt{DriveSending}); still
        accepted today
  \item \textbf{Assessment:} window 0 closed upstream; the residual difference
        is low. Not verified: nothing was executed.
\end{itemize}
\end{frame}

% ============================================================
\section{Finding D}

\begin{frame}{Finding D: Closing a Session}
\begin{center}
\Large Finding D\\[4pt]
\normalsize how a session is closed, and what ends an idle one\\[4pt]
\small 4.19.4.10, 4.19.5 (BTP) \quad and \quad 4.20.3.10, 4.20.4 (PAFTP)
\end{center}
\end{frame}

% ------------------------------------------------------------
\begin{frame}{D --- What the Spec Says About Closing}
\small
\textbf{4.19.4.10 (p.~235), the complete BTP procedure:}
\begin{itemize}
  \item \spec{``To close a BTP session, a GATT client SHALL unsubscribe from
        characteristic C2.''}
  \item \spec{``The GATT server SHALL take this action to indicate closure of
        any BTP session open to the client.''}
  \item \spec{``If a BTP Server needs to close the BTP session, it SHALL
        terminate its BLE connection to the client.''}
\end{itemize}
\textbf{4.20.3.10 (p.~251), PAFTP:} \spec{``A Commissioner MAY terminate the
subscribe instance, to close a PAFTP session. A Commissionable Device MAY
terminate the publish instance, to close a PAFTP session.''}
\vspace{2pt}
\begin{itemize}
  \item Sentence 2 asks the \emph{server} to unsubscribe, but only a GATT
        \emph{client} can subscribe or unsubscribe; sentence 3 gives the server
        a different method
  \item PAFTP: both sides MAY, so no close is required
  \item ``subscribe instance'' / ``publish instance'' appear in 5.4.2.6.3
        (p.~323) and 5.4.3.2 (p.~327) for \emph{creating} one; none defines
        terminating one (defers to WFA-USD, tag M-4)
\end{itemize}
\end{frame}

% ------------------------------------------------------------
\begin{frame}{D2 --- There Is No Close Message}
\small
\textbf{Spec:} one management opcode in each transport, \texttt{0x6C} Handshake
(Table 30, 4.19.3, p.~225; Table 39, 4.20.3, p.~242). No close or reset opcode.\\[4pt]
\textbf{Code:} 13 call sites of \texttt{DoClose} in \texttt{BLEEndPoint.cpp}, 10 in
\texttt{WiFiPAFEndPoint.cpp}. None is triggered by a peer-sent close frame:
\begin{center}
\footnotesize
\begin{tabular}{p{0.20\textwidth}p{0.33\textwidth}p{0.33\textwidth}}
\toprule
\textbf{Trigger} & \textbf{BTP} & \textbf{PAFTP} \\
\midrule
Application & \texttt{Abort}, \texttt{Close} & shutdown paths in \texttt{WiFiPAFLayer.cpp} \\
Link event & peer unsubscribed (\texttt{BleLayer.cpp:583}), connection error (\texttt{:615}) & not traced at the layer level \\
Timer & connect, receive-connection, ack-received, send-ack & the same, plus wait-resource \\
Error & \texttt{Receive()}, \texttt{Send()}, GATT confirmations & \texttt{RxPacketProcess()}, \texttt{Send()}, confirmations \\
\bottomrule
\end{tabular}
\end{center}
\textbf{Assessment:} a structural property. Closing is always a link-layer act
(unsubscribe, disconnect, instance termination) that cannot be authenticated at
this layer. Informational, not an exploit. No claim is made about the secure
session above.
\end{frame}

% ------------------------------------------------------------
\begin{frame}{D1 --- The Idle Timeout in the Spec}
\small
\textbf{4.19.5, Table 35 (pp.~238--239):}
\begin{itemize}
  \item \spec{``BTP\_CONN\_IDLE\_TIMEOUT: The maximum amount of time no unique
        data has been sent over a BTP session before the Central Device must
        close the BTP session. 30 seconds''}
\end{itemize}
\textbf{4.20.4, Table 42 (pp.~253--254):} same for PAFTP, \spec{``\ldots\ before
the Commissioner must close the PAFTP session. 30 seconds''}
\vspace{4pt}
\begin{itemize}
  \item The duty to enforce it is placed on the \emph{connecting} side (Central
        / Commissioner)
  \item The device being paired has no rule for dropping an idle connecting peer
\end{itemize}
\end{frame}

% ------------------------------------------------------------
\begin{frame}[fragile]{D1 --- SDK: the Idle Timeout Is Not Implemented}
\begin{lstlisting}[style=cppsmall]
$ grep -rn "BTP_CONN_IDLE_TIMEOUT" src/            ->  no matches

// src/wifipaf/WiFiPAFConfig.h:144
#define PAFTP_CONN_IDLE_TIMEOUT_MS 30000
// its only use: src/wifipaf/WiFiPAFEndPoint.cpp:76  (a retry count, not a session deadline)
#define WIFIPAF_MAX_RESOURCE_BLOCK_COUNT (PAFTP_CONN_IDLE_TIMEOUT_MS / WIFIPAF_WAIT_RES_TIMEOUT_MS)
\end{lstlisting}
\begin{itemize}\small
  \item \texttt{BLE\_ERROR\_KEEP\_ALIVE\_TIMED\_OUT} and
        \texttt{WIFIPAF\_ERROR\_KEEP\_ALIVE\_TIMED\_OUT} are defined but never
        raised (\texttt{git grep}: only definitions, error-string tables, tests)
  \item By reading: acknowledgements keep re-arming the ack timers
        (\texttt{BTP\_ACK\_TIMEOUT\_MS} = 15\,s, \texttt{BleConfig.h:226}), so a peer
        that \emph{stops} answering is closed
  \item A peer that stays responsive but never does anything is not closed by
        the transport
\end{itemize}
\end{frame}

% ------------------------------------------------------------
\begin{frame}[fragile]{D1 --- BTP and PAFTP Differ at Connection Complete}
\begin{lstlisting}[style=cppsmall]
// src/wifipaf/WiFiPAFEndPoint.cpp:161-163  (HandleConnectComplete)
// Arm ongoing idle session supervision to ensure upper-layer commissioning
// (Phase 2 PASE handshake) initiates within a reasonable deadline.
LogErrorOnFailure(StartAckReceivedTimer());
\end{lstlisting}
\begin{itemize}\small
  \item BTP: \texttt{HandleConnectComplete} (\texttt{BLEEndPoint.cpp:154--176}) and
        \texttt{HandleReceiveConnectionComplete} (\texttt{:178--200}) only cancel
        timers; none is armed. \texttt{git log -S StartAckReceivedTimer} finds
        one BTP commit: the 2020 import.
  \item PAFTP: added by \texttt{650c569666} (2026-06-16, \#72424): \spec{``Resolve
        endpoint pool exhaustion DoS via missing timer \ldots''}
  \item By reading, the first acknowledgement cancels that timer, so it bounds
        only a peer that never acknowledges
\end{itemize}
\begin{center}
\footnotesize
\begin{tabular}{lll}
\toprule
 & \textbf{BTP} & \textbf{PAFTP} \\
\midrule
Endpoint pool (default) & 1 (\texttt{BleConfig.h:78}) & 2 (\texttt{WiFiPAFConfig.h:39}) \\
\bottomrule
\end{tabular}
\end{center}
\end{frame}

% ------------------------------------------------------------
\begin{frame}[fragile]{D1 --- What Does Bound It: the Commissioning Window}
\begin{lstlisting}[style=cppsmall]
// src/app/server/CommissioningWindowManager.cpp:427-438  (abridged)
void CommissioningWindowManager::CloseCommissioningWindow()
{   if (IsCommissioningWindowOpen()) {
        if (mListeningForPASE) {
            // We never established PASE, so never armed a fail-safe and hence
            // can't rely on it expiring to close our BLE connection.  Do that
            // manually here.
            mServer->GetBleLayerObject()->CloseAllBleConnections();
\end{lstlisting}
\begin{itemize}\small
  \item Window length (\texttt{MaxCommissioningTimeout}, \texttt{:318--330}):
        \textbf{15 minutes}, or \textbf{48 hours} on an uncommissioned device
        built with \texttt{CHIP\_DEVICE\_CONFIG\_EXT\_ADVERTISING}
  \item Added in 2022 by \texttt{c916c11c59}: ``Close BLE connections if
        commissioning window closes before PASE establishment''
  \item The comment acknowledges that the transport and the fail-safe do not
        reap the connection; it is in a different layer from the spec section
  \item The PASE and fail-safe timers close the connection too, but are never
        armed if the peer never starts PASE
\end{itemize}
\end{frame}

% ------------------------------------------------------------
\begin{frame}{D1 --- Security Severity}
\textbf{What an attacker needs}
\begin{itemize}\small
  \item Connect to a device in pairing mode, complete the handshake, keep
        acknowledging. No pairing code is needed (before PASE).
\end{itemize}
\textbf{What it gives}
\begin{itemize}\small
  \item By reading: with a pool of 1, one such peer occupies the only BTP
        endpoint until the window closes (15 minutes, up to 48 hours on an
        uncommissioned device with extended advertising); other commissioners
        are refused (endpoint pool full)
  \item No authentication bypass, no data exposure
\end{itemize}
\textbf{Assessment:} low severity (denial of service, owner-opened and
time-bounded window).\\[2pt]
\small Not verified: nothing was executed; platform Bluetooth code was checked by
grep, not read line by line; pool size can differ per platform build.
\end{frame}

% ------------------------------------------------------------
\begin{frame}{D --- Smaller Observation: the Management Bit}
\small
\begin{itemize}
  \item The frame format has an M bit that announces a Management Opcode field
        (4.19.2.1, 4.20.2.1), but only \texttt{0x6C} is defined and nothing says
        what to do with another value
  \item \textbf{BTP:} no M or H bit in \texttt{HeaderFlags}; an M-bit frame
        would have its opcode byte read as an ack or sequence number
  \item \textbf{PAFTP:} \texttt{GetPktSn} skips the opcode byte
        (\texttt{WiFiPAFEndPoint.cpp:859}) but
        \texttt{WiFiPAFTP::HandleCharacteristicReceived}
        (\texttt{WiFiPAFTP.cpp:246--273}) does not, so the two parsers disagree
        on where the sequence number is
\end{itemize}
\textbf{Assessment:} candidate, low severity, medium confidence. Reachable only
by the connected peer; consequence not traced.
\end{frame}

% ============================================================
\section{Summary}

\begin{frame}{Summary}
\footnotesize
\begin{center}
\begin{tabular}{p{0.15\textwidth}p{0.27\textwidth}p{0.22\textwidth}p{0.14\textwidth}}
\toprule
\textbf{Finding} & \textbf{Spec gap} & \textbf{BTP vs PAFTP} & \textbf{Severity} \\
\midrule
C & no rule for a handshake on an established session & closes vs drops; no reason recorded & low \\
Window & no minimum window size & $>0$ vs $\geq 3$; three fixes in 2026 & fixed for 0; residual low \\
D text & server cannot unsubscribe; PAFTP close is optional & --- & low (text) \\
D2 & no close message exists & same in both & informational \\
D1 & idle timeout assigned to the connecting side, not implemented & PAFTP arms a partial timer, BTP none & low \\
\bottomrule
\end{tabular}
\end{center}
\end{frame}

% ------------------------------------------------------------
\begin{frame}{What We Did Not Verify}
\begin{itemize}
  \item Nothing was executed; all code claims are from reading the SDK at
        \texttt{8b3bd4fcc4} (2026-10-08)
  \item The pipeline cannot read figures; the CSA test plans are private
  \item Spec versions after R1.6 and errata were not checked
  \item Platform-specific Bluetooth code was checked by grep, not exhaustively
  \item Over-the-air feasibility of any scenario was not assessed
\end{itemize}
\end{frame}

\end{document}
